% =============================================================================
% Préambule du rapport. Compatible pdfLaTeX, XeLaTeX, LuaLaTeX et Tectonic.
% =============================================================================
\usepackage{iftex}
\ifPDFTeX
  \usepackage[T1]{fontenc}
  \usepackage[utf8]{inputenc}
  \usepackage{lmodern}
  \newcommand{\sansfont}{\sffamily}
\else
  \usepackage{fontspec}
  \newfontfamily\sansfont{DejaVu Sans}
\fi
\usepackage[french]{babel}
\usepackage[a4paper,margin=2.3cm]{geometry}
\usepackage{graphicx}
\usepackage[table]{xcolor}
\usepackage{amssymb}
\usepackage{booktabs}
\usepackage{tabularx}
\usepackage{longtable}
\usepackage{array}
\usepackage{float}
\usepackage{caption}
\usepackage{listings}
\usepackage{tikz}
\usetikzlibrary{positioning,arrows.meta,fit,backgrounds}
\usepackage{fancyhdr}
\usepackage{titlesec}
\usepackage{etoolbox}

% --- Chapitres compacts : pas de saut de page ni de grand titre ----------------
\makeatletter
\patchcmd{\chapter}{\if@openright\cleardoublepage\else\clearpage\fi}{\par\bigskip}{}{}
\makeatother
\titleformat{\chapter}[hang]{\Large\bfseries\color{bleuibam}}{\thechapter.}{0.6em}{}
\titlespacing*{\chapter}{0pt}{1.2em}{0.8em}
\usepackage[hidelinks]{hyperref}

\definecolor{bleuibam}{HTML}{1F4E79}
\definecolor{gristexte}{HTML}{555555}
\definecolor{fondcode}{HTML}{F5F5F5}
\definecolor{orangeafaire}{HTML}{B45309}
\definecolor{bluedark}{RGB}{31,56,100}

% --- Tableaux : bordures complètes, en-tête bleu foncé (\rowcolor{bluedark}) ----
\renewcommand{\arraystretch}{1.25}
\captionsetup[table]{position=top,skip=4pt}
\AtBeginDocument{\renewcommand{\tablename}{Tableau}\renewcommand{\listtablename}{Liste des tableaux}}

% --- En-têtes et pieds de page -------------------------------------------------
\pagestyle{fancy}
\fancyhf{}
\fancyhead[L]{\small\color{gristexte}Groupe 5 SIE, thème 6}
\fancyhead[R]{\small\color{gristexte}API REST de transactions}
\fancyfoot[C]{\thepage}
\renewcommand{\headrulewidth}{0.4pt}

% --- Extraits de terminal et de code -------------------------------------------
\lstset{
  basicstyle=\ttfamily\footnotesize,
  backgroundcolor=\color{fondcode},
  frame=single, rulecolor=\color{gray!40},
  breaklines=true, columns=fullflexible, keepspaces=true,
  showstringspaces=false, upquote=true,
  literate={é}{{\'e}}1 {è}{{\`e}}1 {ê}{{\^e}}1 {à}{{\`a}}1 {ç}{{\c{c}}}1
           {ù}{{\`u}}1 {ô}{{\^o}}1 {î}{{\^i}}1 {É}{{\'E}}1 {«}{{\guillemotleft}}1
           {»}{{\guillemotright}}1 {→}{{$\rightarrow$}}1
}
\renewcommand{\lstlistingname}{Extrait}

% --- \capture{fichier}{légende}{étiquette} -------------------------------------
% Insère une capture de figures/. Tant que le fichier n'existe pas, un cadre
% indique la capture à faire : le rapport compile à tout moment.
\newcommand{\capture}[3]{%
  \begin{figure}[H]
    \centering
    \IfFileExists{figures/#1}{%
      \includegraphics[width=0.9\linewidth,height=0.42\textheight,keepaspectratio]{figures/#1}%
    }{%
      \fbox{\parbox[c][2.6cm][c]{0.85\linewidth}{\centering\color{orangeafaire}
        Capture à insérer : \texttt{\detokenize{figures/#1}}\\[2pt]
        \small\color{gristexte}(liste des captures : figures/LISEZMOI.txt)}}%
    }
    \caption{#2}
    \label{#3}
  \end{figure}}

% --- \afaire{texte} : passage à compléter par l'équipe --------------------------
\newcommand{\afaire}[1]{{\color{orangeafaire}\textbf{[À compléter]} #1}}

\setlength{\parskip}{0.4em}
\setlength{\emergencystretch}{2.5em}
